\begin{theorem}[Equal-characteristic quadratic break parity]\label{mod:cases-wildquadratic-equalcharacteristicbreak}
Let $K/F$ be a finite free Galois valuative extension of nonarchimedean
local fields.  Assume
\[
 \operatorname{char}F=2,\qquad [K:F]=2,\qquad f(K/F)=1.
\]
Thus $e(K/F)=2$, so the extension is totally ramified.  For $t\in\mathbf N$
assume, in the project's integer lower-numbering convention,
\[
 G_t=\operatorname{Gal}(K/F),\qquad G_{t+1}=\{1\}.
\]
Then $t$ is odd.
\LeanFile{LanglandsFirstMainLemma/Cases/WildQuadratic/EqualCharacteristicBreak.lean}
\lean{LanglandsFirstMainLemma.quadraticBreak_odd_equalCharacteristic}
\leanok
\SourceText{Lemma lem:AS-break-parity.}
\uses{mod:ramification-lowergroups,mod:finitefield-artinschreier}
\FormalizationContract
\end{theorem}
\begin{proof}
There is a unique nonidentity $\sigma\in\operatorname{Gal}(K/F)$.  In
characteristic two, fixed-field descent produces $y\in K$ and $a\in F$ with
the Artin--Schreier orientation
\[
 \sigma(y)=y+1,\qquad y^2-y=a.
\]
Since $G_t$ is the full group, $\sigma\in G_0$.  If $y$ were integral, the
defining $G_0$ congruence applied to $y$ would say that $1=\sigma(y)-y$ has
positive order.  Hence $\operatorname{ord}_K(y)<0$.  Comparing the two
unequal orders in $y^2-y=a$, and using $e(K/F)=2$, gives an integer $r>0$
such that
\[
 \operatorname{ord}_K(y)=\operatorname{ord}_F(a)=-r.
\]
Choose such an Artin--Schreier presentation with $r$ minimal; this is the
reduced presentation used below.

Suppose $r=2q$ were even.  Choose a uniformizer $\pi$ of $F$ and put
$u=a\pi^r$.  Then $u$ has order zero and nonzero leading residue
$\bar u$.  On the degree $-2q$ graded piece the leading map induced by
$c\mapsto c^2-c$ is Frobenius.  The finite-field perfectness proved in the
Artin--Schreier node supplies the unique $\bar d$ with
$\bar d^2=\bar u$.  Lift $\bar d$ to a unit $d\in\mathcal O_F$ and set
\[
 c=d/\pi^q,\qquad y'=y-c,
 \qquad a'=a-(c^2-c).
\]
The equality of leading residues gives
\[
 \operatorname{ord}_F(a-c^2)\geq 1-r.
\]
Also $\operatorname{ord}_F(c)=-q\geq1-r$, whence
\[
 \operatorname{ord}_F(a')\geq1-r>-r.
\]
Characteristic two is used again to verify exactly
$(y')^2-y'=a'$ and $\sigma(y')=y'+1$.  The inertia argument makes the new
parameter's order a strictly negative integer, so it yields a positive pole
strictly smaller than $r$, contradicting minimality.  Therefore $r$ is odd.

It remains to identify this pole with the lower break rather than merely
prove parity of an auxiliary parameter.  The pair $(1,y)$ is an $F$-basis
of $K$.  For an integral element $x=A+By$, one has
$\sigma(x)-x=B$.  The two summands have orders $2\operatorname{ord}_F(A)$
and $2\operatorname{ord}_F(B)-r$; their opposite parities prevent
cancellation.  Integrality therefore implies
$r+1\leq\operatorname{ord}_K(B)$, proving $\sigma\in G_r$ directly from the
definition of the lower groups.

Write $r=2q+1$.  For a base uniformizer $\pi$, the element
\[
 \pi_K=\pi^{q+1}y
\]
has order one and is a uniformizer of $K$, while
\[
 \operatorname{ord}_K\bigl(\sigma(\pi_K)-\pi_K\bigr)=r+1.
\]
The necessary uniformizer congruence for lower-group membership then shows
$\sigma\notin G_{r+1}$.  Comparing
$\sigma\in G_r\setminus G_{r+1}$ with
$G_t=\operatorname{Gal}(K/F)$ and $G_{t+1}=\{1\}$ through the decreasing
filtration gives $t=r$.  Hence the project's lower break $t$ is odd.
\end{proof}
